\item Un electrón de cantidad de movimiento $p$ está a una distancia $r$ de un protón estacionario. La energía cinética del electrón es $K = p^2 / 2m_e$, la energía potencial es $U = -k_e e^2 / r$ y la energía total es $E = K + U$. Tratando el átomo como un sistema unidimensional y aplicando la relación de incertidumbre $\Delta p \Delta x \approx \hbar$ con $\Delta x \approx r$:
\begin{enumerate}
    \item estime la incertidumbre de la cantidad de movimiento del electrón en función de $r$,
    \item estime la energía cinética del electrón y la energía total en términos de $r$,
    \item determine el valor de $r$ que minimiza la energía total,
    \item calcule la energía total resultante y
    \item compare su respuesta con las predicciones del modelo de Bohr.
\end{enumerate}